\documentclass[11pt,a4paper]{article}
\usepackage[utf8]{inputenc}
\usepackage[T1]{fontenc}
\usepackage[margin=2.5cm]{geometry}
\usepackage{graphicx,booktabs,amsmath,amssymb,hyperref,xcolor,longtable,listings}
\graphicspath{{../figures/}}
\hypersetup{colorlinks=true,linkcolor=blue!50!black,urlcolor=blue!50!black}
\lstset{basicstyle=\ttfamily\small,breaklines=true,frame=single}

\title{Global Superstore: Retail Analytics \& Profitability Forecasting\\
\large Applied Data Science Project with Python --- MMSE, UAB 2026--2027}
\author{Alejandro Velardo \and Member 2 \and Member 3 \and Member 4}
\date{\today}

\begin{document}
\maketitle
\begin{abstract}
% 150--200 words: problem, data, approach, best model vs. baseline (numbers), key insight on discounts.
TODO
\end{abstract}
\tableofcontents
\newpage

% ---------------------------------------------------------------------------
\section{Introduction \& Problem Statement}
% Context of the dataset; problem categorisation (time-series forecasting / supervised regression);
% research questions & hypotheses (e.g. H1: ML with lag+calendar features beats seasonal naive;
% H2: discount effect on profit is non-linear and category-dependent); team goals.

% ---------------------------------------------------------------------------
\section{Data Cleaning \& Feature Engineering}
\subsection{Data audit}           % schema, types, missingness table (from audit()), duplicates
\subsection{Cleaning decisions}   % each decision + justification + impact (rows affected)
\subsection{Outliers}             % flag vs cap vs drop; why
\subsection{Encoding \& scaling}  % one-hot, cyclical month, StandardScaler for linear models only
\subsection{Temporal aggregation \& features} % panel construction, lags, rolling, discount interactions, leakage control

% ---------------------------------------------------------------------------
\section{Exploratory Data Analysis \& Visual Storytelling}
% Trend & seasonality, discount vs margin, market/category heterogeneity. Each figure: insight + implication.
% \begin{figure}[h]\centering\includegraphics[width=.9\linewidth]{sales_monthly_by_category.pdf}
% \caption{...}\end{figure}

% ---------------------------------------------------------------------------
\section{Theoretical Foundations \& Model Justification}
\subsection{Baselines}  % naive, seasonal naive, moving average
\subsection{Regularised linear models} % Ridge/Lasso objective, bias-variance, log-target, multicollinearity of lags
\subsection{Tree ensembles}            % RF (bagging, variance reduction), gradient boosting (additive stagewise)
\subsection{Why these models for this data} % small n of periods, non-linear discount effects, interactions

% ---------------------------------------------------------------------------
\section{Model Implementation \& Performance Evaluation}
\subsection{Validation strategy} % expanding-window CV, horizon, gap; why not K-fold
\subsection{Metrics}             % WAPE, MAE, RMSE, MASE; why
\subsection{Hyperparameter tuning}
\subsection{Results}             % table baseline vs advanced (mean ± sd over folds), residual analysis, SHAP

% ---------------------------------------------------------------------------
\section{Challenges, Enablers \& Lessons Learned}

% ---------------------------------------------------------------------------
\section{Individual Contribution Breakdown}
\begin{tabular}{@{}lp{10cm}@{}}
\toprule
Member & Role and tasks \\ \midrule
Alejandro Velardo & TODO \\
Member 2 & TODO \\
Member 3 & TODO \\
Member 4 & TODO \\ \bottomrule
\end{tabular}

% ---------------------------------------------------------------------------
\section{Transparent Generative AI Declaration}
See Appendix~\ref{app:ai}.

\appendix
\section{Declaration of Generative AI Tool Usage}\label{app:ai}
% Required: tools used, key prompts, scope of assistance, and >=1 example where the AI output was
% wrong/suboptimal and how the team detected and fixed it. Source: docs/ai_usage_log.md

\end{document}
